% GENERATED FILE. Edit canonical lessons, then run npm run generate:books.
% canonical-source-hash: fnv1a64:44b816ca0042fcb7
% canonical-lessons: MR-C87-mirchi, MR-C87-loni, MR-C87-madh, MR-C87-tel, MR-C87-drakse

\chapter{Chilli, Butter, Honey, Oil, Grapes}
\label{ch:mr-chilli-butter-honey-oil-grapes}
\hlchaptermodality{87}

\begin{chapteropening}
\textbf{By the end of this chapter:} \emph{I can say a chilli, butter, honey, oil and grapes.}

Complete the last lesson of chapter 87: I can say a chilli, butter, honey, oil and grapes.
\end{chapteropening}

\section[\texorpdfstring{\mr{मिरची} (mirchī)}{मिरची (mirchī)}]{\mr{मिरची} (mirchī) — a chilli}

\label{lesson:MR-C87-mirchi}



\begin{quote}
Before the new one: say the Marathi for sugar, then the Marathi for salt.
\end{quote}

\subsection*{You'll want to know: \mr{मिरची}}
\textbf{\mr{मिरची}} — \emph{mirchī} — \textquotedblleft{}a chilli\textquotedblright{}.

\begin{cousinweb}
\textbf{\mr{मिरची}} (\emph{mirchī}) goes back, like Marathi \emph{mirī}, \textquotedblleft{}a peppercorn\textquotedblright{}, to Sanskrit \emph{marīca}, \textquotedblleft{}peppercorn\textquotedblright{} (CDIAL 9875). CDIAL notes the view that the Sanskrit word comes from Austroasiatic. The chilli came later, from the Americas, and took over the pepper's name.
\end{cousinweb}

\begin{grammarlens}[title={What you've built}]
One more word for this chapter.
\end{grammarlens}

\subsection*{Your turn}
\begin{itemize}
\raggedright
  \item \emph{Say it:} \emph{mirchī}
  \item \emph{Say it:} \emph{mirchī}, once more
  \item \emph{Say it:} say \emph{mirchī}
  \item \emph{Recall:} say the Marathi for sugar, then the Marathi for salt, then say \emph{mirchī} again
\end{itemize}

\begin{tcolorbox}[breakable,colback=teal!4,colframe=teal!35!black,title={Before you move on}]
What does \mr{मिरची} mean? (\textquotedblleft{}A chilli\textquotedblright{}.) Say it once more.
\end{tcolorbox}
\section[\texorpdfstring{\mr{लोणी} (loṇī)}{लोणी (loṇī)}]{\mr{लोणी} (loṇī) — butter}

\label{lesson:MR-C87-loni}



\begin{quote}
Before the new one: say the Marathi for salt, then the Marathi for a chilli.
\end{quote}

\subsection*{You'll want to know: \mr{लोणी}}
\textbf{\mr{लोणी}} — \emph{loṇī} — \textquotedblleft{}butter\textquotedblright{}.

\begin{cousinweb}
\textbf{\mr{लोणी}} (\emph{loṇī}) is a tadbhava, a word inherited from Sanskrit and reshaped by sound change. It comes from \emph{navanīta}, \textquotedblleft{}fresh butter\textquotedblright{} (CDIAL 7003). It is built on \emph{nava}, \textquotedblleft{}new\textquotedblright{}.
\end{cousinweb}

\begin{grammarlens}[title={What you've built}]
One more word for this chapter.
\end{grammarlens}

\subsection*{Your turn}
\begin{itemize}
\raggedright
  \item \emph{Say it:} \emph{loṇī}
  \item \emph{Say it:} \emph{loṇī}, once more
  \item \emph{Say it:} say \emph{loṇī}
  \item \emph{Recall:} say the Marathi for salt, then the Marathi for a chilli, then say \emph{loṇī} again
\end{itemize}

\begin{tcolorbox}[breakable,colback=teal!4,colframe=teal!35!black,title={Before you move on}]
What does \mr{लोणी} mean? (\textquotedblleft{}Butter\textquotedblright{}.) Say it once more.
\end{tcolorbox}
\section[\texorpdfstring{\mr{मध} (madh)}{मध (madh)}]{\mr{मध} (madh) — honey}

\label{lesson:MR-C87-madh}



\begin{quote}
Before the new one: say the Marathi for a chilli, then the Marathi for butter.
\end{quote}

\subsection*{You'll want to know: \mr{मध}}
\textbf{\mr{मध}} — \emph{madh} — \textquotedblleft{}honey\textquotedblright{}.

\begin{cousinweb}
\textbf{\mr{मध}} (\emph{madh}) is Sanskrit \emph{madhu}, \textquotedblleft{}honey, mead\textquotedblright{}, lightly reshaped. CDIAL lists the inherited Marathi form as \emph{moh}, \textquotedblleft{}honeycomb, hive\textquotedblright{} (CDIAL 9784).
\end{cousinweb}

\begin{grammarlens}[title={What you've built}]
One more word for this chapter.
\end{grammarlens}

\subsection*{Your turn}
\begin{itemize}
\raggedright
  \item \emph{Say it:} \emph{madh}
  \item \emph{Say it:} \emph{madh}, once more
  \item \emph{Say it:} say \emph{madh}
  \item \emph{Recall:} say the Marathi for a chilli, then the Marathi for butter, then say \emph{madh} again
\end{itemize}

\begin{tcolorbox}[breakable,colback=teal!4,colframe=teal!35!black,title={Before you move on}]
What does \mr{मध} mean? (\textquotedblleft{}Honey\textquotedblright{}.) Say it once more.
\end{tcolorbox}
\section[\texorpdfstring{\mr{तेल} (tel)}{तेल (tel)}]{\mr{तेल} (tel) — oil}

\label{lesson:MR-C87-tel}



\begin{quote}
Before the new one: say the Marathi for butter, then the Marathi for honey.
\end{quote}

\subsection*{You'll want to know: \mr{तेल}}
\textbf{\mr{तेल}} — \emph{tel} — \textquotedblleft{}oil\textquotedblright{}.

\begin{cousinweb}
\textbf{\mr{तेल}} (\emph{tel}) is a tadbhava, a word inherited from Sanskrit and reshaped by sound change. It comes from \emph{taila}, \textquotedblleft{}sesame oil, oil\textquotedblright{} (CDIAL 5958). It is built on \emph{tila}, \textquotedblleft{}sesame\textquotedblright{}.
\end{cousinweb}

\begin{grammarlens}[title={What you've built}]
One more word for this chapter.
\end{grammarlens}

\subsection*{Your turn}
\begin{itemize}
\raggedright
  \item \emph{Say it:} \emph{tel}
  \item \emph{Say it:} \emph{tel}, once more
  \item \emph{Say it:} say \emph{tel}
  \item \emph{Recall:} say the Marathi for butter, then the Marathi for honey, then say \emph{tel} again
\end{itemize}

\begin{tcolorbox}[breakable,colback=teal!4,colframe=teal!35!black,title={Before you move on}]
What does \mr{तेल} mean? (\textquotedblleft{}Oil\textquotedblright{}.) Say it once more.
\end{tcolorbox}
\section[\texorpdfstring{\mr{द्राक्षे} (drākṣe)}{द्राक्षे (drākṣe)}]{\mr{द्राक्षे} (drākṣe) — grapes}

\label{lesson:MR-C87-drakse}



\begin{quote}
Before the new one: say the Marathi for honey, then the Marathi for oil.
\end{quote}

\subsection*{You'll want to know: \mr{द्राक्षे}}
\textbf{\mr{द्राक्षे}} — \emph{drākṣe} — \textquotedblleft{}grapes\textquotedblright{}.

\begin{cousinweb}
\textbf{\mr{द्राक्षे}} (\emph{drākṣe}) is Sanskrit \emph{drākṣā}, \textquotedblleft{}vine, grape\textquotedblright{}, taken over whole as a learned loan, a tatsama (CDIAL 6628).
\end{cousinweb}

\begin{grammarlens}[title={What you've built}]
One more word for this chapter.
\end{grammarlens}

\subsection*{Your turn}
\begin{itemize}
\raggedright
  \item \emph{Say it:} \emph{drākṣe}
  \item \emph{Say it:} \emph{drākṣe}, once more
  \item \emph{Say it:} say \emph{drākṣe}
  \item \emph{Recall:} say the Marathi for honey, then the Marathi for oil, then say \emph{drākṣe} again
\end{itemize}

\begin{tcolorbox}[breakable,colback=teal!4,colframe=teal!35!black,title={Before you move on}]
What does \mr{द्राक्षे} mean? (\textquotedblleft{}Grapes\textquotedblright{}.) Say it once more.
\end{tcolorbox}
